\section{Methods}
\label{sec:methods}

\subsection{Dataset}

We use the Paderborn University KAt bearing dataset \citep{lessmeier2016,katdatacenter}, which contains measurements
of deep-groove ball bearings (type 6203) in an electromechanical drive test rig. Its distinguishing feature for this study is that
it contains both \emph{artificially damaged} bearings (damage introduced by electric discharge machining, drilling or
electric engraving) and bearings with \emph{real damage} that developed during accelerated lifetime tests. We use the
vibration signal (sampled at 64\,kHz). The main analysis uses the operating condition \texttt{N15\_M07\_F10}:
1{,}500\,rpm shaft speed, 0.7\,Nm load torque and 1{,}000\,N radial force. To test whether the findings depend on
this choice, we repeat the evaluation at the dataset's three other conditions: lower load torque (0.1\,Nm), lower
radial force (400\,N) and lower speed (900\,rpm), with the other two parameters unchanged
(Section~\ref{sec:conditions}). Each bearing has 20 recordings of 4\,s per condition.

We consider three classes---healthy, inner-race fault and outer-race fault---and use 29 bearings: 6 healthy, 12 with
artificial damage (7 outer race, 5 inner race) and 11 with real damage (5 outer race, 6 inner race): healthy
K001--K006; artificial outer race KA01, KA03, KA05--KA09; artificial inner race KI01, KI03, KI05, KI07, KI08; real
outer race KA04, KA15, KA16, KA22, KA30; real inner race KI04, KI14, KI16--KI18, KI21. The real damage is fatigue
pitting in nine bearings and plastic deformation (indentations) in KA15 and KA30, with a documented extent of level 1
except for KA16 and KI18 (level 2) and KI16 (level 3) \citep{lessmeier2016}. Bearings with combined inner-
and outer-race damage are excluded. Each recording is cut into non-overlapping 1\,s windows (64{,}000 samples). A
few recordings are slightly shorter than 4\,s and give three windows, so each bearing contributes 76--80 windows
(2{,}309 at the main condition, 2{,}305--2{,}308 at the others).

\subsection{Features}
\label{sec:features}

\paragraph{Fault-frequency features (main setting).} For each window we compute the envelope spectrum
\citep{randall2011}: a fourth-order Butterworth band-pass filter (2--10\,kHz) applied forwards and backwards, the
magnitude of the analytic signal (Hilbert transform) with its mean removed, and the magnitude of its discrete
Fourier transform. For the 6203 bearing ($n = 8$ balls, ball diameter $d = 6.75$\,mm, pitch diameter
$D = 28.55$\,mm, contact angle $0^\circ$) at shaft frequency $f_r = 25$\,Hz, the ball-pass frequencies of the outer
race (BPFO) and of the inner race (BPFI) are
\begin{equation}
  \begin{aligned}
    \mathrm{BPFO} &= \frac{n}{2}\left(1 - \frac{d}{D}\right) f_r \approx 76.4\,\mathrm{Hz},\\
    \mathrm{BPFI} &= \frac{n}{2}\left(1 + \frac{d}{D}\right) f_r \approx 123.6\,\mathrm{Hz}.
  \end{aligned}
\end{equation}
At 900\,rpm ($f_r = 15$\,Hz) these become 45.8\,Hz and 74.2\,Hz.
For each of BPFO and BPFI and their second harmonics, the feature is the largest envelope-spectrum amplitude within
$\pm 2\,\%$ of the target frequency, which allows for the slip of the rolling elements \citep{randall2011}, divided
by the median envelope-spectrum amplitude between 10 and 500\,Hz (the noise floor). These
four ratios are self-normalised: they compare a peak with the spectrum's own noise floor rather than measuring
absolute vibration level. All four are log-transformed before modelling.

\paragraph{All features (ablation).} The ablation adds five time-domain statistics of the raw window: root mean
square, kurtosis, skewness, crest factor and peak-to-peak amplitude (log-transformed except skewness).

\paragraph{Adaptive band (robustness check).} Because a fixed demodulation band might miss impacts that excite
other resonances, we also select the band per window with the fast kurtogram \citep{antoni2007}, which uses spectral
kurtosis \citep{antoni2006} to find the frequency band, in both centre and width, where a signal is most impulsive.
The kurtosis of the complex envelope is computed for every band of the kurtogram's 1/3-binary grid, with each band's
envelope obtained by band-pass filtering in the frequency domain. Bands are limited to 1--30\,kHz and to widths of
at least 1\,kHz, so that the envelope keeps the second harmonic of the inner-race frequency. The band with the
largest kurtosis is used for demodulation, and the same four fault-frequency features are computed in it.

\paragraph{Order of decisions.} Our first analysis used all nine features. We adopted the fault-frequency features as
the main setting after that analysis showed that the amplitude features generalise poorly to unseen bearings of the
\emph{source} domain (Section~\ref{sec:ablation}). Because this choice was made after seeing results, we report both
settings in full.

\subsection{Bearing-level splits}
\label{sec:splits}

Every bearing belongs to exactly one split, so windows from one bearing never appear in both training and
evaluation. The \emph{source domain} contains the 12 artificially damaged bearings and healthy bearings K001, K002
and K004; the \emph{target domain} contains the 11 real-damage bearings and healthy bearings K003, K005 and K006.

The healthy bearings were assigned with care. Their median kurtosis differs strongly: 13.6--15.8 for K001, K002,
K003 and K006, but 4.2--5.1 for K004 and K005. The high values of the first four suggest an impulsive component
unrelated to bearing damage. We place one low-kurtosis healthy bearing on each side, so that the healthy class itself does not shift between domains.
This assignment was made after inspecting the feature distributions but before any model was trained, and the same
split is used for every operating condition.

From the source domain, one bearing per class is held out as the \emph{calibration set}; the remaining source
bearings form the \emph{training set}. Because results depend strongly on which bearings are held out, we repeat the
whole experiment for every possible choice: $3 \times 5 \times 7 = 105$ rotations (healthy $\times$ inner race
$\times$ outer race). Training sets contain 954--957 windows and calibration sets 237--240. The \emph{test set}---all
14 target-domain bearings, 1{,}115 windows---is the same in every rotation.

\subsection{Classifiers}

We compare three classifiers widely used in condition monitoring, implemented with scikit-learn
\citep{pedregosa2011} and the XGBoost library \citep{chen2016}:
\begin{itemize}[nosep]
  \item \textbf{Random Forest} \citep{breiman2001}: 500 trees, minimum 2 samples per leaf, class-balanced weights.
  \item \textbf{Support vector machine (SVM)} \citep{cortes1995}: radial basis function (RBF) kernel, $C = 1$, $\gamma$ = \texttt{scale},
        standardised inputs, class-balanced weights; probabilities from scikit-learn's internal cross-validated
        Platt scaling \citep{platt1999}.
  \item \textbf{XGBoost}: 300 trees, maximum depth 4, learning rate 0.1, row and column subsampling 0.8,
        class-balanced sample weights.
\end{itemize}
All models use a fixed random seed (0) and are fitted on the training set only. Their hyperparameters were set once to
common values and not tuned, for the same reason as for the network (Section~\ref{sec:cnn}).

\subsection{1D-CNN baseline}
\label{sec:cnn}

To test whether an end-to-end deep network avoids the problems of hand-crafted features, we add a one-dimensional
convolutional neural network (1D-CNN) trained on the raw vibration signal, following the design of WDCNN, a deep
convolutional network with a wide first-layer kernel \citep{zhang2017}. Its input is a 4{,}096-sample segment (64\,ms), standardised to zero mean and unit variance,
which removes the absolute vibration level. A first convolution with 16 filters of length 64 and stride 16 is
followed by four convolutions with kernel length 3 (32, 64, 64 and 64 filters); each of the five convolutions is
followed by batch normalisation, rectified linear unit (ReLU) activation and max-pooling by a factor of 2. A fully
connected layer of 100 units, also with batch normalisation and ReLU, and a three-unit output layer complete the
network. The network is trained for
20 epochs on random segments of the training windows (eight per window and epoch), with Adam (learning rate
$10^{-3}$, weight decay $10^{-4}$, batch size 256) and class-balanced cross-entropy. A window's class probabilities
are the mean softmax output over its 15 non-overlapping segments. These settings were fixed in advance, because no
bearings were left for tuning them without using the calibration or test bearings. One network is trained per
rotation, on that rotation's training bearings, and one per held-out bearing for the in-domain reference; their
probabilities then pass through exactly the same calibration, conformal and selective-automation procedures as those
of the other classifiers. As a sanity check, the network and the feature-based classifiers are also trained once on a
random 80/20 split of the source bearings' windows, in which every bearing appears in both the training and the test
set.

\paragraph{Domain adaptation (AdaBN).} To test whether unsupervised domain adaptation makes the network's confidence
trustworthy, we apply adaptive batch normalisation \citep[AdaBN;][]{li2018}, which the WDCNN authors used to adapt
to changing loads \citep{zhang2017}. AdaBN keeps all learned weights and replaces the means and variances stored in
all batch-normalisation layers with those of unlabelled target-domain data, so that every layer's inputs are
standardised in the new domain as they were in the source domain. The adaptation is held out by bearing: each target
bearing is predicted by the rotation's network with statistics estimated on the windows of the other 13 target
bearings, so no test bearing contributes to its own adaptation, and no labels are used. Calibration bearings belong to
the source domain and keep the original statistics, so thresholds and calibrators are fitted exactly as without
adaptation.

\subsection{Calibration and conformal prediction}

Post-hoc calibrators are fitted on the calibration set's predicted probabilities:
\begin{itemize}[nosep]
  \item \textbf{Temperature scaling} \citep{guo2017}: probabilities are rescaled as
        $\mathrm{softmax}(\log \hat{p} / T)$, with a single $T$ chosen to minimise the negative log-likelihood.
  \item \textbf{Isotonic regression} \citep{zadrozny2002}: one-vs-rest isotonic regression per class, followed by
        renormalisation.
\end{itemize}
The uncalibrated (``raw'') probabilities are reported as a third option.

\emph{Split conformal prediction} \citep{vovk2005,angelopoulos2023} uses the non-conformity score
$1 - \hat{p}(y)$ of the true class $y$, computed from the uncalibrated probabilities of the calibration set. The
threshold is the $(k+1)$-th smallest of the $n$ calibration scores, with $k = \lceil (n+1)(1-\alpha) \rceil$ and
$\alpha = 0.1$ (nominal coverage 0.90); this is one rank above the standard split-conformal quantile and therefore
slightly conservative. A class enters a test window's prediction set when its score is at or below the threshold, so sets may
be empty. We use a single threshold (\emph{marginal}) and one threshold per class (\emph{class-conditional},
Mondrian).

\subsection{Selective automation}
\label{sec:selective}

A deployed system acts automatically only when its confidence (the largest class probability) reaches a threshold,
and passes the remaining cases to a human \citep{geifman2017}. For an error target $\varepsilon$ (1\,\% or 5\,\%),
we choose, on the calibration set's (calibrated) probabilities, the lowest threshold whose automatically decided
calibration windows have an error rate of at most $\varepsilon$; ties in confidence are kept together. We then apply this threshold to the test set and report the
\emph{automation rate} (share of test windows decided automatically) and the \emph{automated error} (error rate among
them). As a reference, the \emph{best possible automation} is the automation rate obtained when the threshold is
chosen on the test set itself. If no threshold meets the target on the calibration set, nothing is automated; when we
report how often the target is exceeded across rotations, we count only rotations in which at least one test window
is automated.

\subsection{Metrics}

We report accuracy and per-class recall; the expected calibration error (ECE; top-label, 10 equal-width confidence
bins); the area under the risk--coverage curve (AURC), i.e.\ the mean error among the $k$ most confident windows over
all $k$ \citep{geifman2019}; and conformal empirical coverage and mean prediction-set size. Errors on real damage are
divided into missed faults (a damaged bearing predicted healthy), false alarms (a healthy bearing predicted damaged)
and confusions between the two fault types.

\subsection{In-domain reference}

To separate the artificial-to-real shift from the difficulty of generalising to any unseen bearing, we also evaluate
each classifier on the source domain by leave-one-bearing-out. Each of the 15 source bearings is predicted by a model
trained on the other 14, and metrics are computed on the pooled predictions (raw probabilities).

\subsection{Uncertainty of the estimates}
\label{sec:bootstrap}

With only 3--6 test bearings per class, bearing-to-bearing variation dominates the uncertainty. We therefore use a
hierarchical bootstrap with 2{,}000 draws. Each draw picks one of the 105 rotations at random, then resamples the test
bearings with replacement within each class, keeping the number of bearings per class (3 healthy, 6 inner race,
5 outer race). All metrics are recomputed on the windows of the resampled bearings, with that rotation's calibrators
and thresholds, and we report the mean and the
2.5th--97.5th percentile interval.

\subsection{Calibrating on labelled real bearings}
\label{sec:realcal}

To test whether a small amount of labelled target data restores trustworthy confidence, we move $n = 1$ or 2
target-domain bearings per class from the test set into a separate calibration set, using 50 random selections per
$n$, each combined with a randomly chosen rotation. The classifier is always trained on that rotation's source
training set. We compare three calibration sets, all evaluated on the \emph{same remaining test bearings}:
\emph{source} (the rotation's held-out source bearings), \emph{target} (only the $n$ target bearings per class) and
\emph{combined} (both). Calibrators, conformal thresholds and automation thresholds are fitted on each calibration set
in turn. This experiment is run for the main operating condition only.

\subsection{Implementation}

The pipeline is implemented in Python using NumPy 2.2.6, SciPy 1.15.3, pandas 2.3.3, scikit-learn 1.7.2 and XGBoost
3.2.0; the 1D-CNN is implemented in PyTorch 2.11.0 and was trained on a GPU in Google Colab. Code,
configuration and all result tables are available at \url{https://github.com/amritmalla/bearing-uq-artificial-to-real}; every figure and table can be regenerated
from the Paderborn recordings with the scripts listed in the repository README.
